% GENERATED FILE. Edit canonical lessons, then run npm run generate:books.
% canonical-source-hash: fnv1a64:df88e785a69046c3
% canonical-lessons: RU-C382-provalit, RU-C382-peresdavat, RU-C382-peresdacha, RU-C382-podgotovka, RU-C382-gotovitsya

\chapter{Failing and retaking}
\label{ch:ru-failing-and-retaking}
\hlchaptermodality{382}

\begin{chapteropening}
\textbf{By the end of this chapter:} \emph{I can say that someone failed an exam and must retake it, talk about the resit, and talk about preparing for it.}

Complete the last lesson of chapter 382: I can say that someone failed an exam and must retake it, talk about the resit, and talk about preparing for it.
\end{chapteropening}

\section[\texorpdfstring{\ru{провалить} (provalít)}{провалить (provalít)}]{\ru{провалить} (provalít) — to fail (an exam); to ruin}

\label{lesson:RU-C382-provalit}



\begin{quote}
Before the new one: say the Russian for lecture notes, then the Russian for a point.
\end{quote}

\subsection*{You'll want to know: \ru{провалить}}
\textbf{\ru{провалить}} (\emph{provalít}) — \textquotedblleft{}to fail (an exam); to ruin\textquotedblright{}.

\begin{quote}
\ru{Он} \ru{провалил} \ru{экзамен} \ru{по} \ru{математике}.
\end{quote}

\begin{quote}
\emph{On provalíl ekzámen po matemátike.}
\end{quote}

\begin{quote}
He failed his maths exam.
\end{quote}

Perfective; imperfective \emph{\ru{проваливать}}. Colloquial for exams; the neutral phrase is \emph{\ru{не} \ru{сдать} \ru{экзамен}}.

\begin{cousinweb}
\textbf{\ru{провалить}} is built from \emph{\ru{про}-} \textquotedblleft{}through\textquotedblright{} and \emph{\ru{валить}} \textquotedblleft{}to topple, to knock down\textquotedblright{}: to make a thing fall through. \emph{\ru{Провал}} is a failure, and literally a hole in the ground.
\end{cousinweb}

\begin{grammarlens}[title={What you've built}]
One more word for telling, explaining and describing at B1.
\end{grammarlens}

\subsection*{Your turn}
\begin{itemize}
\raggedright
  \item \emph{Say it:} \emph{provalít}
  \item \emph{Say it:} \emph{provalít}, once more
  \item \emph{Say it:} a sentence of your own with \emph{provalít}
  \item \emph{Recall:} say the Russian for lecture notes, then the Russian for a point, then say \emph{provalít} again
\end{itemize}

\begin{tcolorbox}[breakable,colback=teal!4,colframe=teal!35!black,title={Before you move on}]
What does \ru{провалить} mean? (\textquotedblleft{}To fail (an exam); to ruin\textquotedblright{}.) Say it once more.
\end{tcolorbox}
\section[\texorpdfstring{\ru{пересдавать} (peresdavát)}{пересдавать (peresdavát)}]{\ru{пересдавать} (peresdavát) — to retake (an exam)}

\label{lesson:RU-C382-peresdavat}



\begin{quote}
Before the new one: say the Russian for a point, then the Russian for to fail.
\end{quote}

\subsection*{You'll want to know: \ru{пересдавать}}
\textbf{\ru{пересдавать}} (\emph{peresdavát}) — \textquotedblleft{}to retake (an exam)\textquotedblright{}.

\begin{quote}
\ru{Мне} \ru{придётся} \ru{пересдавать} \ru{химию}.
\end{quote}

\begin{quote}
\emph{Mne pridyótsya peresdavát khímiyu.}
\end{quote}

\begin{quote}
I will have to retake chemistry.
\end{quote}

Imperfective; present \emph{\ru{пересдаю}, \ru{пересдаёшь}}; perfective \emph{\ru{пересдать}}.

\begin{cousinweb}
\textbf{\ru{пересдавать}} is built from \emph{\ru{пере}-} \textquotedblleft{}again\textquotedblright{} and \emph{\ru{сдавать}} \textquotedblleft{}to sit an exam, to hand in\textquotedblright{}, from \emph{\ru{давать}} \textquotedblleft{}to give\textquotedblright{}, the ongoing partner of \emph{\ru{дать}}.
\end{cousinweb}

\begin{grammarlens}[title={What you've built}]
One more word for telling, explaining and describing at B1.
\end{grammarlens}

\subsection*{Your turn}
\begin{itemize}
\raggedright
  \item \emph{Say it:} \emph{peresdavát}
  \item \emph{Say it:} \emph{peresdavát}, once more
  \item \emph{Say it:} a sentence of your own with \emph{peresdavát}
  \item \emph{Recall:} say the Russian for a point, then the Russian for to fail, then say \emph{peresdavát} again
\end{itemize}

\begin{tcolorbox}[breakable,colback=teal!4,colframe=teal!35!black,title={Before you move on}]
What does \ru{пересдавать} mean? (\textquotedblleft{}To retake (an exam)\textquotedblright{}.) Say it once more.
\end{tcolorbox}
\section[\texorpdfstring{\ru{пересдача} (peresdácha)}{пересдача (peresdácha)}]{\ru{пересдача} (peresdácha) — a resit, a retake}

\label{lesson:RU-C382-peresdacha}



\begin{quote}
Before the new one: say the Russian for to fail, then the Russian for to retake.
\end{quote}

\subsection*{You'll want to know: \ru{пересдача}}
\textbf{\ru{пересдача}} (\emph{peresdácha}) — \textquotedblleft{}a resit, a retake\textquotedblright{}.

\begin{quote}
\ru{Пересдача} \ru{будет} \ru{в} \ru{сентябре}.
\end{quote}

\begin{quote}
\emph{Peresdácha búdet v sentyabré.}
\end{quote}

\begin{quote}
The resit will be in September.
\end{quote}

Feminine. From \emph{\ru{пересдать}}, to retake.

\begin{cousinweb}
\textbf{\ru{пересдача}} is built on \emph{\ru{пересдать}} \textquotedblleft{}to retake\textquotedblright{}, from \emph{\ru{пере}-} \textquotedblleft{}again\textquotedblright{} and \emph{\ru{сдать}}, on \emph{\ru{дать}} \textquotedblleft{}to give\textquotedblright{}.
\end{cousinweb}

\begin{grammarlens}[title={What you've built}]
One more word for telling, explaining and describing at B1.
\end{grammarlens}

\subsection*{Your turn}
\begin{itemize}
\raggedright
  \item \emph{Say it:} \emph{peresdácha}
  \item \emph{Say it:} \emph{peresdácha}, once more
  \item \emph{Say it:} a sentence of your own with \emph{peresdácha}
  \item \emph{Recall:} say the Russian for to fail, then the Russian for to retake, then say \emph{peresdácha} again
\end{itemize}

\begin{tcolorbox}[breakable,colback=teal!4,colframe=teal!35!black,title={Before you move on}]
What does \ru{пересдача} mean? (\textquotedblleft{}A resit, a retake\textquotedblright{}.) Say it once more.
\end{tcolorbox}
\section[\texorpdfstring{\ru{подготовка} (podgotóvka)}{подготовка (podgotóvka)}]{\ru{подготовка} (podgotóvka) — preparation}

\label{lesson:RU-C382-podgotovka}



\begin{quote}
Before the new one: say the Russian for to retake, then the Russian for a resit.
\end{quote}

\subsection*{You'll want to know: \ru{подготовка}}
\textbf{\ru{подготовка}} (\emph{podgotóvka}) — \textquotedblleft{}preparation\textquotedblright{}.

\begin{quote}
\ru{Подготовка} \ru{к} \ru{экзамену} \ru{занимает} \ru{много} \ru{времени}.
\end{quote}

\begin{quote}
\emph{Podgotóvka k ekzámenu zanimáyet mnógo vrémeni.}
\end{quote}

\begin{quote}
Preparing for an exam takes a lot of time.
\end{quote}

Feminine. \emph{\ru{Подготовка} \ru{к}} + dative.

\begin{cousinweb}
\textbf{\ru{подготовка}} is built on \emph{\ru{подготовить}} \textquotedblleft{}to prepare\textquotedblright{}, from \emph{\ru{готовый}} \textquotedblleft{}ready\textquotedblright{}, inherited from Proto-Slavic \emph{*gotovŭ}.
\end{cousinweb}

\begin{grammarlens}[title={What you've built}]
One more word for telling, explaining and describing at B1.
\end{grammarlens}

\subsection*{Your turn}
\begin{itemize}
\raggedright
  \item \emph{Say it:} \emph{podgotóvka}
  \item \emph{Say it:} \emph{podgotóvka}, once more
  \item \emph{Say it:} a sentence of your own with \emph{podgotóvka}
  \item \emph{Recall:} say the Russian for to retake, then the Russian for a resit, then say \emph{podgotóvka} again
\end{itemize}

\begin{tcolorbox}[breakable,colback=teal!4,colframe=teal!35!black,title={Before you move on}]
What does \ru{подготовка} mean? (\textquotedblleft{}Preparation\textquotedblright{}.) Say it once more.
\end{tcolorbox}
\section[\texorpdfstring{\ru{готовиться} (gotóvitsya)}{готовиться (gotóvitsya)}]{\ru{готовиться} (gotóvitsya) — to prepare oneself, to study for}

\label{lesson:RU-C382-gotovitsya}



\begin{quote}
Before the new one: say the Russian for a resit, then the Russian for preparation.
\end{quote}

\subsection*{You'll want to know: \ru{готовиться}}
\textbf{\ru{готовиться}} (\emph{gotóvitsya}) — \textquotedblleft{}to prepare oneself, to study for\textquotedblright{}.

\begin{quote}
\ru{Я} \ru{готовлюсь} \ru{к} \ru{экзамену}.
\end{quote}

\begin{quote}
\emph{Ya gotóvlyus k ekzámenu.}
\end{quote}

\begin{quote}
I am preparing for the exam.
\end{quote}

Imperfective, with the present \emph{\ru{готовлюсь}, \ru{готовишься}}; perfective \emph{\ru{подготовиться}}. \emph{\ru{Готовиться} \ru{к}} + dative. Without \emph{-\ru{ся}}, \emph{\ru{готовить}} is to cook or to prepare something.

\begin{cousinweb}
\textbf{\ru{готовиться}} is the reflexive of \emph{\ru{готовить}}, from \emph{\ru{готовый}} \textquotedblleft{}ready\textquotedblright{}, inherited from Proto-Slavic \emph{*gotovŭ}: to make yourself ready.
\end{cousinweb}

\begin{grammarlens}[title={What you've built}]
One more word for telling, explaining and describing at B1.
\end{grammarlens}

\subsection*{Your turn}
\begin{itemize}
\raggedright
  \item \emph{Say it:} \emph{gotóvitsya}
  \item \emph{Say it:} \emph{gotóvitsya}, once more
  \item \emph{Say it:} a sentence of your own with \emph{gotóvitsya}
  \item \emph{Recall:} say the Russian for a resit, then the Russian for preparation, then say \emph{gotóvitsya} again
\end{itemize}

\begin{tcolorbox}[breakable,colback=teal!4,colframe=teal!35!black,title={Before you move on}]
What does \ru{готовиться} mean? (\textquotedblleft{}To prepare oneself, to study for\textquotedblright{}.) Say it once more.
\end{tcolorbox}
